\subsection*{Central simple algebras and 2-cocycles}

Fix a finite Galois extension $L$ of $k$, and let $G=\operatorname{Gal}(L/k)$. Define $\mathcal{A}(L/k)$ to be the class of central simple algebras $A$ over $k$ containing $L$ and of degree $[A:k]=[L:k]^2$ (hence, $L$ equals its centralizer in $A$).

Fix an $A\in\mathcal{A}(L/k)$. For every $\sigma\in G$, Corollary 2.11 of the Noether-Skolem theorem shows that there exists an element $e_\sigma\in A$ such that
\[
\sigma a=e_\sigma a e_\sigma^{-1}\quad\text{for all } a\in L. \tag{38}
\]
Moreover, $e_\sigma$ is determined by (38) up to multiplication by an element of $L^\times$, because if $f_\sigma$ has the same property, then $f_\sigma^{-1}e_\sigma$ centralizes $L$. Note that (38) can be written as
\[
e_\sigma\cdot a=\sigma a\cdot e_\sigma\quad\text{for all } a\in L, \tag{39}
\]
which says that moving $e_\sigma$ past $a\in L$ replaces it with $\sigma a$. Clearly $e_\sigma e_\tau$ has the property (38) for $\sigma\tau$, and so
\[
e_\sigma e_\tau=\varphi(\sigma,\tau)e_{\sigma\tau} \tag{40}
\]
for some $\varphi(\sigma,\tau)\in L^\times$. Note that
\[
e_\rho(e_\sigma e_\tau)=e_\rho(\varphi(\sigma,\tau)e_{\sigma\tau})=\rho\varphi(\sigma,\tau)\cdot\varphi(\rho,\sigma\tau)\cdot e_{\rho\sigma\tau}
\]
and
\[
(e_\rho e_\sigma)e_\tau=\varphi(\rho,\sigma)e_{\rho\sigma}e_\tau=\varphi(\rho,\sigma)\varphi(\rho\sigma,\tau)\cdot e_{\rho\sigma\tau}.
\]
Therefore the associative law implies that $\varphi$ is a 2-cocycle. It is even a normalized 2-cocycle if we choose $e_1=1$. A different choice of $e_\sigma$'s leads to a cohomologous 2-cocycle, and so we have a well-defined map $A\mapsto\gamma(A)\colon\mathcal{A}(L/k)\to H^2(L/k)$.

\noindent\textsc{Theorem 3.11}\quad The map $\gamma\colon\mathcal{A}(L/k)\to H^2(L/k)$ is surjective, and its fibres are the isomorphism classes.

We first need a lemma.

\noindent\textsc{Lemma 3.12}\quad Let $A\in\mathcal{A}(L/k)$, and define $e_\sigma$ to satisfy (38). Then the set $(e_\sigma)_{\sigma\in G}$ is a basis for $A$ as a left vector space over $L$.

\noindent\textsc{Proof.}\quad Note that
\[
\dim_L(A)=\frac{\dim_k(A)}{\dim_k(L)}=n,
\]
and so it suffices to show that the $e_\sigma$ are linearly independent. Suppose not, and let $(e_\sigma)_{\sigma\in J}$ be a maximal linearly independent set. If $\tau\notin J$, then
\[
e_\tau=\sum a_\sigma e_\sigma
\]
for some $a_\sigma\in L$. Let $a\in L$. When we compute $e_\tau a$ in two different ways,
\[
e_\tau a=\tau a\cdot e_\tau=\sum_{\sigma\in J}\tau a\cdot a_\sigma e_\sigma,
\]
\[
e_\tau a=\sum_{\sigma\in J}a_\sigma e_\sigma\cdot a=\sum_{\sigma\in J}a_\sigma\cdot\sigma a\cdot e_\sigma
\]
we find that $\tau a\cdot a_\sigma=\sigma a\cdot a_\sigma$ for all $\sigma\in J$. For at least one $\sigma\in J$, $a_\sigma\neq0$, and then the equation shows that $\tau=\sigma$, contradicting the fact that $\tau\notin J$. Therefore $J=G$. \hfill$\square$

Now $A$ is uniquely determined by the following properties: $A\supset L$; $(e_\sigma)_{\sigma\in G}$ is a basis for $A$ as an $L$-vector space; multiplication in $A$ satisfies the equation (38) and (40).

Let $A'\in\mathcal{A}(L/k)$ and suppose that $\gamma(A)=\gamma(A')$. The condition implies that we can choose bases $(e_\sigma)$ and $(e_\sigma')$ for $A$ and $A'$ satisfying (38) and (40) with the same 2-cocycle $\varphi$. The map $\sum a_\sigma e_\sigma\mapsto\sum a_\sigma e_\sigma'\colon A\to A'$ is an isomorphism of $k$-algebras.

Next suppose that $A$ and $A'$ are isomorphic elements of $\mathcal{A}(L/k)$. The Noether-Skolem theorem allows us to choose the isomorphism $f\colon A\to A'$ so that $f(L)=L$ and $f|_L$ is the identity map. If $e_\sigma$ satisfies condition (40) for $A$, then $f(e_\sigma)$ satisfies (40) for $A'$. With the choices $(e_\sigma)$ and $(f(e_\sigma))$, $A$ and $A'$ define the same cocycle.

These remarks show that the map $A\mapsto\gamma(A)$ defines a injection
\[
\mathcal{A}(L/k)/\approx\to H^2(L/k).
\]
To show that the map is surjective, we construct an inverse.

Let $\varphi\colon G\times G\to L^\times$ be a normalized 2-cocycle. Define $A(\varphi)$ to be the $L$-vector space with basis $(e_\sigma)_{\sigma\in G}$ endowed with the multiplication given by (38) and (40). Then $e_1$ is an identity element for multiplication, and the cocycle condition (exactly) shows that
\[
e_\rho(e_\sigma e_\tau)=(e_\rho e_\sigma)e_\tau.
\]
It follows that $A(\varphi)$ is a $k$-algebra. We identify $L$ with the subfield $Le_1$ of $A(\varphi)$.

\noindent\textsc{Lemma 3.13}\quad The algebra $A(\varphi)$ is central simple over $K$.

\noindent\textsc{Proof.}\quad Let $\alpha=\sum a_\sigma e_\sigma$ centralize $L$, and let $a\in L$. On comparing $a\alpha=\sum aa_\sigma\cdot e_\sigma$ with $\alpha a=\sum a_\sigma(\sigma a)\cdot e_\sigma$, we find that $a_\sigma=0$ for $\sigma\neq1$, and so $\alpha=a_1e_1\in L$. Therefore, the centralizer of $L$ in $A(\varphi)$ is $L$.

Let $\alpha$ lie in the centre of $A(\varphi)$. Then $\alpha$ centralizes $L$, and so $\alpha\in L$, say $\alpha=ae_1$, $a\in L$. On comparing $e_\sigma\cdot\alpha=(\sigma a)e_\sigma$ with $\alpha\cdot e_\sigma=ae_\sigma$, we see that $\alpha\in k$. Thus $A(\varphi)$ is central.

Let $\mathfrak{A}$ be a two-sided ideal in $A(\varphi)$; in particular, $\mathfrak{A}$ is an $L$-subspace of $A(\varphi)$. If $\mathfrak{A}$ contains one element $e_\sigma$, then (40) shows that it contains all, and so equals $A(\varphi)$. Suppose $\mathfrak{A}\neq0$, and let $\alpha=\sum a_\sigma e_\sigma$ be a primordial element of $\mathfrak{A}$, with say $a_{\sigma_0}=1$. If $a_{\sigma_1}\neq0$, $\sigma_1\neq\sigma_0$, then for every $a\in L$,
\[
(\sigma_1 a)\cdot\alpha-\alpha\cdot a=\sum a_\sigma(\sigma_1 a-\sigma a)e_\sigma\in\mathfrak{A}.
\]
If $a$ is chosen so that $\sigma_1 a\neq\sigma_0 a$, then this element is nonzero but has fewer nonzero coefficients than $\alpha$, contradicting its primordality. Therefore, $\alpha=e_{\sigma_0}$, and we have shown that $\mathfrak{A}=A(\varphi)$. \hfill$\square$

Let $\varphi$ and $\varphi'$ be cohomologous 2-cocycles, say,
\[
a(\sigma)\cdot\sigma a(\tau)\cdot\varphi'(\sigma,\tau)=a(\sigma\tau)\cdot\varphi(\sigma,\tau)
\]
for some map $a\colon G\to L^\times$. One checks immediately that the $L$-linear map $A(\varphi)\to A(\varphi')$ sending $e_\sigma$ to $a(\sigma)e_\sigma'$ is an isomorphism of $k$-algebras. Therefore $\varphi\mapsto A(\varphi)$ defines a map $H^2(L/K)\to\mathcal{A}(L/k)/\approx$, which is clearly inverse to $A\mapsto\gamma(A)$. This completes the proof of Theorem 3.11.

The algebras $A(\varphi)$ are called \textbf{\textit{crossed-product algebra}}. Before group cohomology existed, 2-cocycles $\varphi\colon G\times G\to L^\times$ were called \textbf{\textit{factor sets}} (and sometimes still are).

To show that this map is bijective, it suffices (after Theorem 3.11) to show that the map $A\mapsto[A]\colon\mathcal{A}(L/k)/\approx\to\operatorname{Br}(L/k)$ is bijective.

If $A$ and $A'$ are similar central simple algebras over $k$, then (see 2.13) there exists a central division algebra $D$ such that $A\sim D\sim A'$, say, $A\approx M_n(D)$, $A'\approx M_{n'}(D)$. But if $[A:k]=[A':k]$, then $n=n'$, and so $A\approx A'$. This proves that the map $\mathcal{A}(L/k)/\approx\to\operatorname{Br}(L/k)$ is injective, and 3.6 proves that it is surjective.

\noindent\textsc{Lemma 3.15}\quad For any two 2-cocycles $\varphi$ and $\varphi'$, $A(\varphi+\varphi')\sim A(\varphi)\otimes_k A(\varphi')$.

\noindent\textsc{Proof.}\quad The proof is a little messy because we have to recognize $A(\varphi)\otimes_k A(\varphi')$, not as a crossed-product algebra, but as matrix algebra over a crossed-product algebra. I merely sketch the proof.\footnote[4]{See Blanchard 1972, pp.\ 94--95, or Farb and Dennis 1993, p.\ 126--128, for the details.}

Set $A=A(\varphi)$, $B=A(\varphi')$, and $C=A(\varphi+\varphi')$. Regard $A$ and $B$ as left $L$-modules (using left multiplication), and define
\[
V=A\otimes_L B.
\]
Concretely, $V$ is the largest quotient space of $A\otimes_k B$ such that
\[
la\otimes_L b=a\otimes_L lb
\]
holds for all $a\in A$, $b\in B$, $l\in L$.

The $k$-vector space $V$ has a unique right $A\otimes_k B$-module structure such that
\[
(a'\otimes_L b')(a\otimes_k b)=a'a\otimes_L b'b\,,\quad\text{all } a',a\in A\,,\ b',b\in B,
\]
and a unique left $C$-module structure such that
\[
(le_\sigma'')(a\otimes_L b)=le_\sigma a\otimes_L e_\sigma' b\,,\quad\text{all } l\in L,\ \sigma\in G,\ a\in A,\ b\in B.
\]
Here $(e_\sigma)$, $(e_\sigma')$, and $(e_\sigma'')$ are the standard bases for $A=A(\varphi)$, $B=A(\varphi')$, and $C=A(\varphi+\varphi')$ respectively.

The two actions commute, and so the right action of $A\otimes_k B$ on $V$ defines a homomorphism of $k$-algebras
\[
f\colon(A\otimes_k B)^{\mathrm{opp}}\to\operatorname{End}_C(V).
\]
This homomorphism is injective because $A\otimes_k B$ (and hence its opposite) is simple. Since both $(A\otimes_k B)^{\mathrm{opp}}$ and $\operatorname{End}_C(V)$ have degree $n^4$ over $k$, where $n=[L:k]$, $f$ is an isomorphism. As we showed in (1.21), any two modules over a simple ring of the same $k$-dimension are isomorphic, and it follows that $V\approx C^n$ as a $C$-module. Hence
\[
\operatorname{End}_C(V)\approx\operatorname{End}_C(C^n)=M_n(C^{\mathrm{opp}}),
\]
and on composing this isomorphism with $f$ we obtain an isomorphism of $k$-algebras
\[
(A\otimes_k B)^{\mathrm{opp}}\to M_n(C)^{\mathrm{opp}}.
\]
The same map can be interpreted as an isomorphism
\[
A\otimes_k B\to M_n(C). \hfill\square
\]
